\chapter{\abstractname}

3D Gaussian Splatting (3DGS) has emerged as a fast, high-quality alternative
to Neural Radiance Fields for novel view synthesis, but inserting new content
into a 3DGS scene remains an open problem: the implicit, per-Gaussian
representation does not natively expose surface geometry, material, or
lighting parameters needed for plausible compositing.

This thesis presents \emph{InsertGS}, a pipeline for inserting mesh-based
objects into rooms reconstructed as 3D Gaussian Splatting scenes. The system
combines (i)~a GPU rasterizer that converts a captured 3DGS room into a
cubemap environment map suitable for image-based lighting, (ii)~a physically
based Blender stage that renders the inserted object under that environment,
and (iii)~a compositing pass that places the rendered object back into novel
3DGS views with consistent shading and shadows. A Gradio dashboard exposes
each stage so practitioners can run individual operations or full
end-to-end batches, and a lightweight web viewer based on SuperSplat lets
users inspect input scenes and intermediate results in the browser.

We evaluate the pipeline on a curated set of indoor 3DGS reconstructions and
several inserted objects, showing that the cubemap-based IBL stage produces
visually consistent lighting compared to ad-hoc compositing, and that the
batch CLI scales to many cases without manual intervention.
